% =====================================================================
% Tesis 2 -- Propuesta de proyecto
% Deteccion de anti-spoofing de voz para espanol multi-acento:
% aumentacion del corpus MARSA y entrenamiento de detectores de linea base
%
% Documento AUTOCONTENIDO (bibliografia embebida). Compila en Overleaf
% con pdflatex directamente. Idioma: espanol.
%
% Estructura solicitada:
%   1. Objetivo general del proyecto
%   2. Antecedentes
%   3. Justificacion
%   4. Objetivos especificos del proyecto
%   5. Ambiente de implementacion y restricciones consideradas
%   6. Etapas de desarrollo
%   7. Cronograma de actividades
%   8. Resultados esperados
%   9. Bibliografia
%
% Marcado en rojo con \PENDIENTE{...} lo que aun no esta definido.
% =====================================================================
\documentclass[11pt]{article}

\usepackage[utf8]{inputenc}
\usepackage[T1]{fontenc}
\usepackage[spanish,es-noquoting]{babel}
\usepackage{lmodern}
\usepackage[margin=1in]{geometry}
\usepackage{booktabs}
\usepackage{graphicx}
\usepackage{amsmath,amssymb}
\usepackage{array}
\usepackage{longtable}
\usepackage{xcolor}
\usepackage{enumitem}
\usepackage[colorlinks=true,linkcolor=black,citecolor=black,urlcolor=blue]{hyperref}
\usepackage[numbers,sort&compress]{natbib}

\newcommand{\PENDIENTE}[1]{\textcolor{red}{\textbf{[PENDIENTE: #1]}}}

\title{\textbf{Detecci\'on de \emph{anti-spoofing} de voz para espa\~nol
multi-acento: aumentaci\'on del corpus MARSA y entrenamiento de
detectores de l\'inea base con evaluaci\'on de generalizaci\'on
\emph{cross-corpus}}\\[6pt]
\large Propuesta de Tesis 2 (continuaci\'on)}

\author{%
Tom\'as Acosta Bernal \\
\small Asesor: Rub\'en Manrique \\
\small Universidad de los Andes, Bogot\'a, Colombia
}
\date{}

\begin{document}
\maketitle

% =====================================================================
\section{Objetivo general del proyecto}

Dise\~nar, entrenar y evaluar un sistema de detecci\'on de
\emph{anti-spoofing} de voz para espa\~nol multi-acento que sea robusto a
condiciones reales de canal y a manipulaciones parciales de la se\~nal,
utilizando el corpus \textbf{MARSA} (\emph{Multi-accent Audio Recordings
for Spanish Anti-spoofing}) construido en la Tesis~1. El proyecto
extiende el corpus con una capa de aumentaci\'on ac\'ustica y una capa de
aumentaci\'on aplicada al \emph{partial spoofing}, entrena tres
detectores de \emph{anti-spoofing} del estado del arte sobre el corpus
resultante, y cuantifica su capacidad de generalizaci\'on evaluando en un
corpus de espa\~nol independiente (HISPASpoof). El objetivo \'ultimo es
producir la evidencia de validaci\'on t\'ecnica (\emph{baseline} de EER)
que permita completar y someter el art\'iculo de datos (\emph{Data
Descriptor}) de MARSA a la revista \emph{Scientific Data}.

% =====================================================================
\section{Antecedentes}

\subsection{El problema del \emph{anti-spoofing} de voz}

Los sistemas de verificaci\'on autom\'atica de locutor (ASV) se despliegan
a gran escala en banca, telefon\'ia y control de acceso. Su principal
amenaza de seguridad es el \emph{spoofing} de voz: un atacante sintetiza
o clona la voz de un objetivo para enga\~nar al verificador. La serie de
retos ASVspoof estableci\'o los conjuntos de datos y protocolos de
evaluaci\'on est\'andar para este problema
\citep{asvspoof2019,asvspoof2021}, y desde entonces se han desarrollado
numerosas contramedidas neuronales \citep{rawnet2,aasist,tak2022ssl}. Un
hallazgo recurrente y bien documentado es que estas contramedidas
generalizan mal: un detector casi perfecto sobre su distribuci\'on de
entrenamiento se degrada de forma severa ante sistemas de s\'intesis,
condiciones de grabaci\'on o corpus no vistos
\citep{muller2022generalize,muller2024harder}.

\subsection{La ausencia del espa\~nol y el rol de la Tesis 1}

Los conjuntos de datos que anclan esta investigaci\'on son casi
exclusivamente en ingl\'es o mandar\'in. El espa\~nol (la segunda lengua
materna m\'as hablada del mundo, distribuida en dos continentes y en
variedades ac\'usticamente distintas) ha carecido de un corpus de
\emph{anti-spoofing} a gran escala. Adem\'as, la variaci\'on
inter-dialectal del espa\~nol es suficientemente grande como para
degradar de forma no uniforme el desempe\~no de modelos de voz de \'ultima
generaci\'on
\citep{kantharuban2023dialectgap,nacimiento2024biases,zuluaga2023commonaccent},
por lo que un detector no puede asumirse transferible entre variedades.

La \textbf{Tesis 1} atac\'o esta ausencia construyendo el corpus MARSA. A
partir de 35\,927 grabaciones aut\'enticas (\emph{bonafide}) de 1\,567
locutores en siete acentos (peninsular m\'as seis variedades
latinoamericanas), se generaron:
\begin{itemize}[nosep]
  \item \textbf{186\,477 ataques de \emph{full spoof}} mediante seis
  sistemas de clonaci\'on de voz TTS de c\'odigo abierto (Fish-Speech
  \citep{fishspeech}, Qwen3-TTS \citep{qwen3tts}, OpenVoice
  \citep{openvoice}, Chatterbox \citep{chatterbox}, OuteTTS
  \citep{outetts} y OmniVoice \citep{omnivoice});
  \item \textbf{86\,766 ataques de \emph{partial spoof}}, en los que se
  reemplaza de una a tres palabras de una grabaci\'on aut\'entica por
  habla clonada, mediante un motor de \emph{splicing} sensible a la
  energ\'ia con siete t\'ecnicas de fundido y una variante de
  homogeneizaci\'on de fronteras (\emph{boundary jitter}) que elimina un
  atajo conocido de los detectores \citep{negroni2024splicing}.
\end{itemize}
El total base del corpus es de aproximadamente 309\,170 enunciados. El
\emph{partial spoofing} es la amenaza m\'as dif\'icil de detectar, pues la
mayor parte de la se\~nal es aut\'entica \citep{zhang2023partialspoof}; el
\'unico recurso multiling\"ue previo que incluye espa\~nol, HQ-MPSD
\citep{li2025hqmpsd}, aporta del orden de $1.3\times10^{4}$ muestras de un
solo sistema TTS y un solo acento, sin aumentaci\'on. MARSA supera ese
recurso en escala, diversidad de sistemas, cobertura de acentos y capas
de aumentaci\'on.

\subsection{Lo que falta: entrenar y evaluar detectores sobre MARSA}

La Tesis 1 entreg\'o el \emph{corpus}, pero no un \emph{detector}
entrenado ni evaluado sobre \'el. Sin un resultado de detecci\'on
(\emph{baseline} de EER) el corpus no est\'a validado como recurso \'util,
y la secci\'on de \emph{Technical Validation} del art\'iculo de datos
queda incompleta. Adem\'as, existen recursos de espa\~nol independientes
(en particular HISPASpoof \citep{hispaspoof}, un corpus de forensia de
habla en espa\~nol con seis acentos) que permiten medir la
generalizaci\'on \emph{cross-corpus} de un detector entrenado en MARSA.
Finalmente, han aparecido detectores de \emph{anti-spoofing} recientes y
potentes que a\'un no se han evaluado sobre habla en espa\~nol:
Nes2Net \citep{nes2net}, una arquitectura anidada ligera sobre
representaciones de modelos fundacionales (WavLM); el sistema RAPTOR /
DF\_Arena \citep{dfarena}, basado en \emph{backbones} auto-supervisados
compactos; y HoliAntiSpoof \citep{holiantispoof}, un modelo de lenguaje
de audio (Audio-LLM) que reformula el an\'alisis de \emph{spoofing} como
una tarea de generaci\'on de texto (detecci\'on, m\'etodo, regi\'on e
impacto sem\'antico). La Tesis 2 se ocupa precisamente de este vac\'io.

% =====================================================================
\section{Justificaci\'on}

\begin{itemize}[leftmargin=1.2em]
  \item \textbf{Relevancia de seguridad.} El fraude por voz sint\'etica
  sobre canales telef\'onicos es una amenaza creciente en Am\'erica
  Latina. Un detector robusto y espec\'ifico para espa\~nol multi-acento
  atiende una necesidad de despliegue real que los detectores entrenados
  en ingl\'es no cubren.

  \item \textbf{Robustez de canal.} Los ataques reales llegan a trav\'es
  de micr\'ofonos, salas reverberantes y c\'odecs de telefon\'ia/VoIP. La
  aumentaci\'on ac\'ustica (reverberaci\'on, ruido y c\'odecs reales,
  aplicados de forma simple y apilada) es necesaria para que el detector
  no colapse fuera del laboratorio.

  \item \textbf{La frontera dif\'icil.} No existe un corpus aumentado de
  \emph{partial spoofing} en espa\~nol. Aumentar el \emph{partial spoof}
  de MARSA y medir su efecto sobre la detecci\'on es una contribuci\'on
  original y directamente pertinente a la amenaza m\'as evasiva.

  \item \textbf{Generalizaci\'on como problema central.} El problema
  abierto del \emph{anti-spoofing} es la generalizaci\'on. Entrenar en
  MARSA y evaluar en HISPASpoof cuantifica de forma honesta cu\'anto
  transfiere un detector de espa\~nol entre corpus, acentos y sistemas de
  s\'intesis.

  \item \textbf{Cierre del art\'iculo de datos.} El \emph{baseline} de EER
  es la \'unica pieza que falta para someter el \emph{Data Descriptor} de
  MARSA a \emph{Scientific Data}; la Tesis 2 la produce.
\end{itemize}

% =====================================================================
\section{Objetivos espec\'ificos del proyecto}

\begin{enumerate}[label=\textbf{OE\arabic*.},leftmargin=2.6em]
  \item \textbf{Consolidar la aumentaci\'on ac\'ustica de MARSA a escala
  de corpus.} Ejecutar y verificar la capa de aumentaci\'on
  (reverberaci\'on + ruido, c\'odecs reales, y el \emph{gate} de apilado
  60/40) sobre el corpus completo, en los factores de multiplicaci\'on
  previstos ($\times$2, $\times$3, $\times$5, $\times$10), con la
  pol\'itica de \emph{loudness} uniforme y el separador de identificador
  de sistema con ``\texttt{|}'' para muestras apiladas.

  \item \textbf{Dise\~nar y aplicar aumentaci\'on al corpus de
  \emph{partial spoofing}.} Extender la capa de aumentaci\'on a las
  muestras de \emph{partial spoof} (incluida la variante \emph{jittered}),
  preservando las etiquetas a nivel de trama, y documentar el efecto
  sobre las m\'etricas de calidad.

  \item \textbf{Preparar el protocolo de entrenamiento y evaluaci\'on.}
  Definir particiones disjuntas por locutor (\emph{train}/\emph{dev}/%
  \emph{eval}), \emph{data loaders} en formato ASVspoof2019 LA, y el mapeo
  de etiquetas necesario para evaluar sobre HISPASpoof.

  \item \textbf{Ajustar (\emph{fine-tuning}) tres detectores del estado
  del arte sobre MARSA.} Nes2Net \citep{nes2net}, RAPTOR/DF\_Arena
  \citep{dfarena} y HoliAntiSpoof \citep{holiantispoof} (este \'ultimo
  v\'ia adaptador LoRA sobre su base de 7B par\'ametros).

  \item \textbf{Evaluar el desempe\~no \emph{in-corpus} de forma
  estratificada.} Reportar EER global y desagregado por (i)~sistema de
  ataque, (ii)~\emph{tier} de \emph{partial spoof} (W1/W2/W3), (iii)~grupo
  de acento, y (iv)~condici\'on de aumentaci\'on.

  \item \textbf{Evaluar la generalizaci\'on \emph{cross-corpus}.} Entrenar
  en MARSA y evaluar en HISPASpoof (y, de ser viable, en direcci\'on
  inversa) para medir la transferencia entre corpus de espa\~nol.

  \item \textbf{Completar y someter el \emph{Data Descriptor} de MARSA.}
  Integrar los resultados de \emph{baseline} en la secci\'on de
  \emph{Technical Validation} y preparar el dep\'osito p\'ublico del
  corpus (Zenodo) para el env\'io a \emph{Scientific Data}.
\end{enumerate}

% =====================================================================
\section{Ambiente de implementaci\'on y restricciones consideradas}

\subsection{Ambiente de implementaci\'on}

\begin{itemize}[nosep]
  \item \textbf{Servidor de c\'omputo:} un servidor con 4$\times$
  NVIDIA A40 (46\,GB de VRAM cada una). Es una m\'aquina
  \emph{compartida} con otros investigadores.
  \item \textbf{Flujo de dos m\'aquinas:} edici\'on de c\'odigo y control
  de versiones en una m\'aquina local; ejecuci\'on, entrenamiento
  e inferencia en el servidor de c\'omputo.
  \item \textbf{Aislamiento de entornos:} cada sistema de ataque y cada
  detector se ejecuta en su propio entorno virtual aislado, para no
  afectar los entornos de otros investigadores.
  \item \textbf{Recursos externos de aumentaci\'on:} conjunto de respuestas
  al impulso RIRS\_NOISES \citep{ko2017rir} y ruido MUSAN \citep{musan};
  c\'odecs reales de telefon\'ia y VoIP.
  \item \textbf{M\'etricas y herramientas:} EER como m\'etrica primaria;
  NISQA \citep{nisqa} y ECAPA-TDNN \citep{ecapa} para calidad y similitud
  de locutor; Parakeet-TDT \citep{tdt,parakeet} para transcripci\'on y
  alineamiento.
\end{itemize}

\subsection{Restricciones consideradas}

\begin{itemize}[leftmargin=1.2em]
  \item \textbf{Tiempo:} el proyecto dispone de un \'unico semestre
  (aproximadamente 16 semanas). El alcance se prioriza en consecuencia.
  \item \textbf{GPU compartida:} se usa \emph{una sola} GPU a la vez, lo
  que limita el paralelismo de entrenamiento.
  \item \textbf{HoliAntiSpoof es un modelo de 7B par\'ametros.} A\'un con
  LoRA, su \emph{fine-tuning} e inferencia son costosos en una A40
  compartida. Se contempla como objetivo de mayor riesgo; si el tiempo o
  el c\'omputo no alcanzan, se degrada a evaluaci\'on \emph{zero-shot}
  (sin \emph{fine-tuning}) y se prioriza Nes2Net y RAPTOR/DF\_Arena.
  \item \textbf{Madurez de los modelos.} Los tres detectores son
  publicaciones muy recientes (\emph{preprints} 2025--2026); existe riesgo
  de reproducibilidad (c\'odigo, pesos, dependencias). Se reserva tiempo
  de integraci\'on por modelo.
  \item \textbf{Licencias.} DF\_Arena/RAPTOR se distribuye bajo licencia
  \emph{no comercial}; su uso es acad\'emico. \PENDIENTE{Verificar la
  licencia exacta de HoliAntiSpoof antes de redistribuir salidas.}
  \item \textbf{Alineaci\'on con HISPASpoof.} La evaluaci\'on
  \emph{cross-corpus} requiere mapear el formato y las etiquetas de
  HISPASpoof al protocolo de MARSA; se considera una etapa expl\'icita.
  \item \textbf{Desbalance de acentos.} Los tres acentos m\'as peque\~nos
  (argentino, peruano, venezolano) suman 6.6\% del corpus; el EER por
  acento para esos grupos tendr\'a mayor varianza y se reporta con esa
  salvedad.
  \item \textbf{Dep\'osito de datos.} El env\'io a \emph{Scientific Data}
  exige dep\'osito p\'ublico del corpus (Zenodo/Figshare); se reserva
  tiempo para la preparaci\'on del dep\'osito.
\end{itemize}

% =====================================================================
\section{Etapas de desarrollo}

\begin{description}[leftmargin=0em,style=nextline]
  \item[Etapa 1. Aumentaci\'on ac\'ustica a escala (OE1).]
  Ejecuci\'on y verificaci\'on de la capa de reverberaci\'on+ruido y de
  c\'odecs reales, con el \emph{gate} de apilado, en los cuatro factores.
  Validaci\'on de que la aumentaci\'on no introduce atajos de etiqueta
  (\emph{loudness} uniforme, fracci\'on limpia por clase).

  \item[Etapa 2. Aumentaci\'on del \emph{partial spoof} (OE2).]
  Extensi\'on de la capa de aumentaci\'on al corpus de \emph{partial
  spoof} preservando etiquetas a nivel de trama; auditor\'ia de calidad.

  \item[Etapa 3. Protocolo de datos (OE3).]
  Particiones disjuntas por locutor; \emph{data loaders} en formato
  ASVspoof2019 LA; mapeo de HISPASpoof al protocolo de MARSA.

  \item[Etapa 4. Integraci\'on de los detectores (OE4).]
  Puesta en marcha de los tres modelos en el servidor (entornos,
  pesos, dependencias); pruebas de humo de inferencia.

  \item[Etapa 5. \emph{Fine-tuning} sobre MARSA (OE4).]
  Ajuste de Nes2Net y RAPTOR/DF\_Arena; ajuste de HoliAntiSpoof v\'ia LoRA
  (o \emph{zero-shot} como contingencia).

  \item[Etapa 6. Evaluaci\'on \emph{in-corpus} estratificada (OE5).]
  C\'alculo de EER global y desagregado por sistema de ataque, \emph{tier},
  acento y condici\'on de aumentaci\'on.

  \item[Etapa 7. Evaluaci\'on \emph{cross-corpus} (OE6).]
  Entrenamiento en MARSA y evaluaci\'on en HISPASpoof; an\'alisis de la
  brecha de generalizaci\'on.

  \item[Etapa 8. Consolidaci\'on y art\'iculo de datos (OE7).]
  Integraci\'on de resultados en \emph{Technical Validation}; preparaci\'on
  del dep\'osito Zenodo; escritura final y env\'io a \emph{Scientific
  Data}; redacci\'on del documento de tesis.
\end{description}

% =====================================================================
\section{Cronograma de actividades}

El cronograma cubre 16 semanas (un semestre). La marca ``$\bullet$''
indica actividad principal en la semana correspondiente.

\begin{center}
\footnotesize
\setlength{\tabcolsep}{3pt}
\renewcommand{\arraystretch}{1.25}
\begin{longtable}{p{5.2cm}*{16}{c}}
\toprule
\textbf{Actividad / Semana} &
1&2&3&4&5&6&7&8&9&10&11&12&13&14&15&16 \\
\midrule
\endhead
E1 Aumentaci\'on ac\'ustica a escala
 &$\bullet$&$\bullet$& & & & & & & & & & & & & & \\
E2 Aumentaci\'on \emph{partial spoof}
 & &$\bullet$&$\bullet$& & & & & & & & & & & & & \\
E3 Protocolo y particiones + mapeo HISPASpoof
 & & &$\bullet$&$\bullet$& & & & & & & & & & & & \\
E4 Integraci\'on de los 3 detectores
 & & & &$\bullet$&$\bullet$& & & & & & & & & & & \\
E5 \emph{Fine-tuning} Nes2Net + RAPTOR/DF\_Arena
 & & & & &$\bullet$&$\bullet$&$\bullet$& & & & & & & & & \\
E5 \emph{Fine-tuning} HoliAntiSpoof (LoRA)
 & & & & & & &$\bullet$&$\bullet$&$\bullet$& & & & & & & \\
E6 Evaluaci\'on \emph{in-corpus} estratificada
 & & & & & & & & &$\bullet$&$\bullet$&$\bullet$& & & & & \\
E7 Evaluaci\'on \emph{cross-corpus} (HISPASpoof)
 & & & & & & & & & & &$\bullet$&$\bullet$& & & & \\
E8 Technical Validation + dep\'osito Zenodo
 & & & & & & & & & & & &$\bullet$&$\bullet$&$\bullet$& & \\
E8 Escritura y env\'io (Scientific Data + tesis)
 & & & & & & & & & & & & & &$\bullet$&$\bullet$&$\bullet$ \\
\bottomrule
\end{longtable}
\end{center}

\noindent\textbf{Contingencia.} Si en la semana~9 el \emph{fine-tuning} de
HoliAntiSpoof no es viable en la GPU compartida, se sustituye por
evaluaci\'on \emph{zero-shot} de ese modelo y se reasigna el tiempo a un
an\'alisis m\'as profundo de los dos detectores restantes. El camino
cr\'itico para el \emph{Data Descriptor} (E6--E8) queda protegido.

% =====================================================================
\section{Resultados esperados}

\begin{enumerate}[label=\textbf{RE\arabic*.},leftmargin=2.6em]
  \item \textbf{Corpus MARSA aumentado} en sus dos capas (ac\'ustica y de
  \emph{partial spoof}), con etiquetas preservadas y documentaci\'on de
  calidad.
  \item \textbf{Tres detectores ajustados} (o, en el peor caso, dos
  ajustados y uno evaluado \emph{zero-shot}) sobre MARSA.
  \item \textbf{Tabla de EER \emph{in-corpus} estratificada} por sistema
  de ataque, \emph{tier} de \emph{partial spoof}, acento y condici\'on de
  aumentaci\'on.
  \item \textbf{Resultados de generalizaci\'on \emph{cross-corpus}}
  (entrenar en MARSA, evaluar en HISPASpoof), con an\'alisis de la brecha
  espa\~nol-vs-espa\~nol entre corpus.
  \item \textbf{Art\'iculo de datos completo} de MARSA, con la secci\'on de
  \emph{Technical Validation} llena, y su \textbf{dep\'osito p\'ublico}
  preparado, sometido a \emph{Scientific Data}.
  \item \textbf{Documento de tesis} que integra la construcci\'on
  (Tesis~1) y la evaluaci\'on (Tesis~2) del corpus.
\end{enumerate}

\noindent\emph{Nota de alcance.} Esta propuesta enumera el conjunto
completo de actividades previstas; la priorizaci\'on por tiempo y por
c\'omputo compartido implica que el n\'ucleo garantizado son OE1--OE6 con
al menos dos detectores, y que OE7 (tercer detector, direcci\'on inversa
de \emph{cross-corpus}) se persigue en la medida en que los recursos lo
permitan.

% =====================================================================
% Bibliografia embebida (autocontenida; no requiere .bib externo).
% =====================================================================
\begin{thebibliography}{99}

\bibitem{asvspoof2019}
M.~Todisco \emph{et al.},
``ASVspoof 2019: Future horizons in spoofed and fake audio detection,''
en \emph{Proc. Interspeech}, 2019, pp.~1008--1012.

\bibitem{asvspoof2021}
J.~Yamagishi \emph{et al.},
``ASVspoof 2021: Accelerating progress in spoofed and deepfake speech
detection,'' en \emph{Proc. ASVspoof 2021 Workshop}, 2021.

\bibitem{rawnet2}
H.~Tak \emph{et al.},
``End-to-end anti-spoofing with RawNet2,'' en \emph{Proc. ICASSP}, 2021.

\bibitem{aasist}
J.-w.~Jung \emph{et al.},
``AASIST: Audio anti-spoofing using integrated spectro-temporal graph
attention networks,'' en \emph{Proc. ICASSP}, 2022.

\bibitem{tak2022ssl}
H.~Tak \emph{et al.},
``Automatic speaker verification spoofing and deepfake detection using
wav2vec 2.0 and data augmentation,'' en \emph{Proc. Odyssey}, 2022,
pp.~112--119.

\bibitem{muller2022generalize}
N.~M.~M\"uller \emph{et al.},
``Does audio deepfake detection generalize?'' en \emph{Proc.
Interspeech}, 2022, pp.~2783--2787.

\bibitem{muller2024harder}
N.~M.~M\"uller \emph{et al.},
``Harder or different? Understanding generalization of audio deepfake
detection,'' en \emph{Proc. Interspeech}, 2024.

\bibitem{habla}
P.~A.~Tamayo-Fl\'orez, R.~Manrique, y B.~Pereira Nunes,
``HABLA: A dataset of Latin American Spanish accents for voice
anti-spoofing,'' en \emph{Proc. Interspeech}, 2023, pp.~1963--1967.

\bibitem{fishspeech}
S.~Liao \emph{et al.},
``Fish-Speech: Leveraging large language models for advanced
multilingual text-to-speech synthesis,'' \emph{arXiv:2411.01156}, 2024.

\bibitem{qwen3tts}
Qwen Team, Alibaba Group,
``Qwen3-TTS technical report,'' \emph{arXiv:2601.15621}, 2026.

\bibitem{openvoice}
Z.~Qin, W.~Zhao, X.~Yu, y X.~Sun,
``OpenVoice: Versatile instant voice cloning,'' \emph{arXiv:2312.01479},
2024.

\bibitem{chatterbox}
Resemble AI,
``Chatterbox: Open-source text-to-speech with neural watermarking,''
publicaci\'on de software, 2025.
\url{https://github.com/resemble-ai/chatterbox}

\bibitem{outetts}
OuteAI,
``OuteTTS: LLM-based text-to-speech via discrete audio tokens,''
publicaci\'on de software, 2025. \url{https://huggingface.co/OuteAI}

\bibitem{omnivoice}
H.~Zhu \emph{et al.},
``OmniVoice: Towards omnilingual zero-shot text-to-speech with diffusion
language models,'' \emph{arXiv:2604.00688}, 2026.

\bibitem{zhang2023partialspoof}
L.~Zhang, X.~Wang, E.~Cooper, N.~Evans, y J.~Yamagishi,
``The PartialSpoof database and countermeasures for the detection of
short fake speech segments embedded in an utterance,''
\emph{IEEE/ACM Trans. Audio, Speech, Lang. Process.}, vol.~31,
pp.~813--825, 2023.

\bibitem{li2025hqmpsd}
M.~Li \emph{et al.},
``HQ-MPSD: A multilingual artifact-controlled benchmark for partial
deepfake speech detection,'' \emph{arXiv:2512.13012}, 2025.

\bibitem{negroni2024splicing}
V.~Negroni, D.~Salvi, P.~Bestagini, y S.~Tubaro,
``Analyzing the impact of splicing artifacts in partially fake speech
signals,'' en \emph{Proc. ASVspoof 5 Workshop}, 2024.

\bibitem{hispaspoof}
M.~Risques, K.~Bhagtani, A.~K.~S.~Yadav, y E.~J.~Delp,
``HISPASpoof: A new dataset for Spanish speech forensics,''
en \emph{Proc. ICASSP}, 2026.

\bibitem{kantharuban2023dialectgap}
A.~Kantharuban, I.~Vuli\'c, y A.~Korhonen,
``Quantifying the dialect gap and its correlates across languages,''
en \emph{Findings of EMNLP}, 2023.

\bibitem{zuluaga2023commonaccent}
J.~Zuluaga-Gomez, S.~Ahmed, D.~Visockas, y C.~Subakan,
``CommonAccent: Exploring large acoustic pretrained models for accent
classification based on Common Voice,''
en \emph{Proc. Interspeech}, 2023, pp.~5291--5295.

\bibitem{nacimiento2024biases}
E.~Nacimiento-Garc\'ia, H.~S.~D\'iaz-Kaas-Nielsen, y
C.~S.~Gonz\'alez-Gonz\'alez,
``Gender and accent biases in AI-based tools for Spanish: A comparative
study between Alexa and Whisper,''
\emph{Applied Sciences}, vol.~14, no.~11, p.~4734, 2024.

\bibitem{nes2net}
T.~Liu, D.-T.~Truong, R.~K.~Das, K.~A.~Lee, y H.~Li,
``Nes2Net: A lightweight nested architecture for foundation model driven
speech anti-spoofing,'' \emph{arXiv:2504.05657}, 2025.

\bibitem{dfarena}
A.~Kulkarni, S.~Dowerah, A.~Kulkarni, T.~Alum\"ae, y M.~Magimai-Doss,
``Do compact SSL backbones matter for audio deepfake detection? A
controlled study with RAPTOR,'' \emph{arXiv:2603.06164}, 2026.
Modelo: \url{https://huggingface.co/Speech-Arena-2025/DF_Arena_1B_V_1}

\bibitem{holiantispoof}
``HoliAntiSpoof: Audio LLM for holistic speech anti-spoofing,''
\emph{arXiv:2602.04535}, 2026.
Modelo: \url{https://huggingface.co/wsntxxn/HoliAntiSpoof}

\bibitem{ecapa}
B.~Desplanques, J.~Thienpondt, y K.~Demuynck,
``ECAPA-TDNN: Emphasized channel attention, propagation and aggregation
in TDNN based speaker verification,''
en \emph{Proc. Interspeech}, 2020, pp.~3830--3834.

\bibitem{nisqa}
G.~Mittag, B.~Naderi, A.~Chehadi, y S.~M\"oller,
``NISQA: A deep CNN-self-attention model for multidimensional speech
quality prediction with crowdsourced datasets,''
en \emph{Proc. Interspeech}, 2021, pp.~2127--2131.

\bibitem{tdt}
H.~Xu, F.~Jia, S.~Majumdar, S.~Watanabe, y B.~Ginsburg,
``Efficient sequence transduction by jointly predicting tokens and
durations,'' en \emph{Proc. ICML}, 2023.

\bibitem{parakeet}
NVIDIA NeMo Team,
``Parakeet-TDT-0.6B-v3: Multilingual token-and-duration transducer ASR,''
publicaci\'on de modelo, 2024.
\url{https://huggingface.co/nvidia/parakeet-tdt-0.6b-v3}

\bibitem{musan}
D.~Snyder, G.~Chen, y D.~Povey,
``MUSAN: A music, speech, and noise corpus,'' \emph{arXiv:1510.08484},
2015.

\bibitem{ko2017rir}
T.~Ko, V.~Peddinti, D.~Povey, M.~L.~Seltzer, y S.~Khudanpur,
``A study on data augmentation of reverberant speech for robust speech
recognition,'' en \emph{Proc. ICASSP}, 2017.

\end{thebibliography}

\end{document}
